%%%%%%%%%%%%%%%%%%%%%%%%%%%%%%%%%%%%%%%%%%%%%%%%%%%%%%%%%%%%%%%%%%%%%%%%%%%%%%%%
%% G2 -- MCM/ICM 2027 SUMMARY SHEET (Problem C: Olympic medal tables).
%%
%% WHAT THIS FILE IS
%%   The summary sheet only.  It is written on the official COMAP template, so
%%   it compiles on its own to exactly ONE page -- the local acceptance test for
%%   a summary sheet.  To submit, paste the block between "Begin Summary" and
%%   "End Summary" below into the paper skeleton, in place of the template's red
%%   sample text, so that this text is page 1 of the submitted PDF.
%%
%% PAGE BUDGET (25 pages)
%%   The 25-page limit covers the Summary Sheet, table of contents, full
%%   solution, references, notes and appendices.  The only excluded section is
%%   the "Report on Use of AI", which must be placed after the 25 pages.
%%
%% THE TWO FILLED FIELDS
%%   \Problem -- the chosen problem letter (C here), NOT the template "ABCDEF".
%%   \Team    -- the Team Control Number.  The value below is a stand-in (no
%%               real control number was available when this sheet was written);
%%               replace 2400001 with the team's actual control number, in both
%%               \Team and the submitted file name "<control number>.pdf".
%%
%% PAPER SIZE
%%   The page size is not set here (the official \geometry only fixes margins);
%%   if your compiler does not produce US Letter or A4, set it explicitly.
%%%%%%%%%%%%%%%%%%%%%%%%%%%%%%%%%%%%%%%%%%%%%%%%%%%%%%%%%%%%%%%%%%%%%%%%%%%%%%%%
%%%%%%%%%%%%%%%%%%%%%%%%%%%%%%%%%%%%%%%%
%% MCM/ICM LaTeX Template %%
%% 2027 MCM/ICM           %%
%%%%%%%%%%%%%%%%%%%%%%%%%%%%%%%%%%%%%%%%
\documentclass[12pt]{article}
\usepackage{geometry}
\geometry{left=1in,right=0.75in,top=1in,bottom=1in}

%%%%%%%%%%%%%%%%%%%%%%%%%%%%%%%%%%%%%%%%
% Replace ABCDEF in the next line with your chosen problem
% and replace 1111111 with your Team Control Number
\newcommand{\Problem}{C}
\newcommand{\Team}{2400001}
%%%%%%%%%%%%%%%%%%%%%%%%%%%%%%%%%%%%%%%%

\usepackage{newtxtext}
\usepackage{amsmath,amssymb,amsthm}
\usepackage{newtxmath} % must come after amsXXX

\usepackage[pdftex]{graphicx}
\usepackage{xcolor}
\usepackage{fancyhdr}
\lhead{Team \Team}
\rhead{}
\cfoot{}

\newtheorem{theorem}{Theorem}
\newtheorem{corollary}[theorem]{Corollary}
\newtheorem{lemma}[theorem]{Lemma}
\newtheorem{definition}{Definition}

%%%%%%%%%%%%%%%%%%%%%%%%%%%%%%%%
\begin{document}
% Place your graphic files in the same directory as your main document
\graphicspath{{.}}
\DeclareGraphicsExtensions{.pdf, .jpg, .tif, .png}
\thispagestyle{empty}
\vspace*{-16ex}
\centerline{\begin{tabular}{*3{c}}
	\parbox[t]{0.3\linewidth}{\begin{center}\textbf{Problem Chosen}\\ \Large \textcolor{red}{\Problem}\end{center}}
	& \parbox[t]{0.3\linewidth}{\begin{center}\textbf{2027\\ MCM/ICM\\ Summary Sheet}\end{center}}
	& \parbox[t]{0.3\linewidth}{\begin{center}\textbf{Team Control Number}\\ \Large \textcolor{red}{\Team}\end{center}}	\\
	\hline
\end{tabular}}
%%%%%%%%%%% Begin Summary %%%%%%%%%%%
\begin{center}
\textbf{2028 Olympic Medal Predictions Based on Random Forest Model}
\end{center}

We predict the medal table event by event instead of extrapolating national totals. From the
1896--2024 athlete, medal-table, host-country and programme data we build a four-dimensional
capability vector for every athlete or team entry -- past total medals, past gold medals, best
past result, and past score per appearance, each shifted to exclude the current Games -- and
train one random forest per sport for the probability of winning a medal and another for the
probability of winning gold; sports with fewer than five records or a single training class fall
back to their historical medal rate. Medals are then allocated within each event (one gold, then
the two highest remaining probabilities) rather than summed, which prevents inflated totals.
Both steps assume the 2024 entrant pool and the 2024 programme carry over to 2028.

A 1{,}000-run Monte Carlo over the fitted probabilities gives 95\% prediction intervals. Total
medals: USA 113--142, China 105--125, Great Britain 52--75, Australia and Italy 45--65, France
43--63, Japan 40--58, Netherlands 31--47, Germany 29--47. Gold medals: USA 37--53, China 31--49.
Against 2024, the United States, China and Italy are our major improvers and France, Australia
and Germany the major decliners. Held-out discrimination is good (ROC AUC 0.974 and 0.983 on two
folds; PR-AUC 0.926 on the best fold, 0.639 on the worst), but the calibration curve
over-predicts above $p = 0.7$, so the intervals above are mildly optimistic.

More than 60 NOCs have never won a medal. Repeated over 500 simulated Games, the same draw gives
an expected 4.78 first-time medallists at Los Angeles 2028, with 4 to 7 new NOCs in a typical
simulation. The strongest individual odds are Samoa (SAM) 0.41, Mali (MLI) 0.27, Guam (GUM) 0.25,
Papua New Guinea (PNG) 0.24 and Vanuatu (VAN) 0.24.

Poisson regression connects programme size to medals. For USA swimming one extra event raises the
expected medal count by 21.1\% ($\beta_1 = 0.1912$, $p < 0.001$); the estimate survives negative
binomial (0.2378) and zero-inflated Poisson (0.1626, lowest AIC 92.33) specifications and
leave-one-year-out refitting (0.18--0.20, except 0.2305 without 1964). China's table tennis is
saturated: $\beta_1 = 0.1273$ with $p = 0.59$, so extra events add nothing to an already complete
dominance. Hosting is worth about 74.8 total medals for the USA ($p < 0.001$) and 24.7 for China
($p = 0.055$), consistent with hosts choosing extra events in their own strengths.

We split a country's winning capability ($0$--$3$ per event) into athlete and coach components.
Lang Ping's USA women's volleyball added 2.00 and 1.33 capability levels in 2008 and 2012 (mean
1.16), and 1.67 for China in 2016: one coach moved a team by one to two medal grades. Applying
that model, we advise India to invest in badminton (worth 1--2 bronze medals) and field hockey,
Sweden in swimming, and Romania in rowing.

Our distinctive finding is distributional. The NOCs still without a medal cluster in Africa, the
Pacific and Latin America, where the binding constraint is programme access rather than athlete
quality; for those countries we recommend adding events in sports the region already contests, as
the cheapest route to a first medal, with the odds quantified above. National committees should
treat programme composition -- worth up to 21\% per event in a strong sport, and far more to a
host -- as a preparation target set years in advance rather than at the Games.
%%%%%%%%%%% End Summary %%%%%%%%%%%

%%%%%%%%%%%%%%%%%%%%%%%%%%%%%%
% Everything below is the template's page-break scaffolding for the FULL paper.
% It is commented out here so that this file compiles to exactly one page.
% Uncomment it (or simply paste the summary block above into the paper skeleton)
% when the summary sheet is used as page 1 of the submitted report.
%\clearpage
%\pagestyle{fancy}
%\tableofcontents   % uncomment to generate a Table of Contents
%\newpage
%\setcounter{page}{1}
%\rhead{Page \thepage\ }
%%%%%%%%%%%%%%%%%%%%%%%%%%%%%%
\end{document}
\end
